\chapterbackmatter{Solutions to Weinberg Exercises}
\ExerciseEditorialNote

\begin{enumerate}
  \WeinbergSolution{1}
  In any representation,
  \[
    [D_\mu,[D_\nu,D_\lambda]]
    +[D_\nu,[D_\lambda,D_\mu]]
    +[D_\lambda,[D_\mu,D_\nu]]=0
  \]
  by the Jacobi identity for operator commutators.  Using
  \(
    [D_\nu,D_\lambda]=-it_\alpha F_{\alpha\nu\lambda}
  \)
  and letting the outer derivative act on an arbitrary matter field gives
  \[
    -it_\alpha\bigl(
      D_\mu F_{\alpha\nu\lambda}
      D_\nu F_{\alpha\lambda\mu}
      D_\lambda F_{\alpha\mu\nu}
    \bigr)=0.
  \]
  Here the derivative on the field strength is the adjoint covariant
  derivative,
  \[
    D_\mu F_{\alpha\nu\lambda}
    =\partial_\mu F_{\alpha\nu\lambda}
      C_{\alpha\beta\gamma}A_{\beta\mu}F_{\gamma\nu\lambda}.
  \]
  Faithfulness is not needed: the same identity follows componentwise by
  inserting the definition of \(F\) and using the Jacobi identity for
  \(C_{\alpha\beta\gamma}\).  This proves the stated Bianchi identity.

  \WeinbergSolution{2}
  An infinitesimal gauge transformation gives
  \[
    \delta f_\alpha
    =\partial_i(D_i\lambda)_\alpha,\qquad
    (D_i\lambda)_\alpha
    =\partial_i\lambda_\alpha
      +C_{\alpha\gamma\beta}A_{\gamma i}\lambda_\beta .
  \]
  Thus the Faddeev--Popov operator and ghost action are
  \[
    \mathcal F_{\alpha x,\beta y}
    =\partial_i^x(D_i^x)_{\alpha\beta}\delta^4(x-y),
    \qquad
    I_{\rm GH}=\int d^4x\,
      \omega_\alpha^*\partial_i(D_i\omega)_\alpha .
  \]
  After a spatial integration by parts,
  \[
    \mathcal L_{\rm GH}
    =-\partial_i\omega_\alpha^*\partial_i\omega_\alpha
      +C_{\alpha\beta\gamma}
       (\partial_i\omega_\alpha^*)A_{\gamma i}\omega_\beta ,
  \]
  where antisymmetry of the structure constants has been used in the
  interaction term.  The field-independent kernel is
  \(\delta_{\alpha\beta}\boldsymbol\nabla^2\).  In the propagator convention
  of Eq.~\eqref{eq:15.6.12}, its inverse is
  \[
    \Delta_{\alpha\beta}(x,y)
    =\delta_{\alpha\beta}\delta(x^0-y^0)
      \int\frac{d^3\mathbf p}{(2\pi)^3}\,
      \frac{-e^{i\mathbf p\cdot(\mathbf x-\mathbf y)}}{\mathbf p^2}.
  \]
  Equivalently its momentum kernel is
  \(-\delta_{\alpha\beta}/\mathbf p^2\).  The overall sign is reversed if
  \(-\boldsymbol\nabla^2\), rather than \(\boldsymbol\nabla^2\), is used as
  the positive spatial operator.  The important feature is the absence of a
  \(p^0\) pole: Coulomb-gauge ghosts propagate instantaneously.

  \WeinbergSolution{3}
  In electrodynamics \(\delta A_\mu=\partial_\mu\lambda\), so
  \[
    \delta f
      =\left(\Box+2cA^\mu\partial_\mu\right)\lambda .
  \]
  The ghost Lagrangian therefore is
  \[
    \mathcal L_{\rm GH}
      =\omega^*\left(\Box+2cA^\mu\partial_\mu\right)\omega
      =-\partial_\mu\omega^*\partial^\mu\omega
        +2c\,\omega^*A^\mu\partial_\mu\omega ,
  \]
  up to a surface term in the second form.  Its free kernel is still
  \(\Box\), and hence
  \[
    \Delta(x,y)
      =\int\frac{d^4p}{(2\pi)^4}\,
        \frac{e^{ip\cdot(x-y)}}{p^2-i\epsilon}.
  \]
  Unlike ghosts in a linear Abelian gauge, these ghosts do not decouple:
  the nonlinear term in \(f\) produces a photon--ghost--antighost vertex.

  \WeinbergSolution{4}
  For the compact Lie algebras relevant here, complexify the algebra and
  choose a Cartan subalgebra of rank \(r\).  Its root decomposition gives
  \[
    \dim\mathfrak g=r+2N_+ ,
  \]
  where \(N_+\) is the number of positive roots.  If
  \(\dim\mathfrak g=4\), then \(r\) is even.  The possibility \(r=4\)
  has no roots and is Abelian, hence is not simple.  The only remaining
  possibility is \(r=2\) and \(N_+=1\).  But then the roots are only
  \(\pm\alpha\), and they span a one-dimensional subspace of the
  two-dimensional dual Cartan algebra.  A nonzero Cartan element \(H\)
  annihilated by \(\alpha\) commutes both with the Cartan subalgebra and with
  the two root spaces.  It is therefore central, again contradicting
  simplicity.  Thus no compact simple Lie algebra has four generators.

  \WeinbergSolution{5}
  Write
  \(A_\mu=t_\alpha A_{\alpha\mu}\).  Successive transport along the sides of
  an infinitesimal parallelogram with side vectors \(a^\mu\) and \(b^\mu\)
  gives, through second order,
  \[
    \psi_{\rm final}-\psi_{\rm initial}
      =i\,t_\alpha F_{\alpha\mu\nu}(X)a^\mu b^\nu\psi(X).
  \]
  For a general small loop the oriented area bivector is
  \[
    \Sigma^{\mu\nu}
      =\frac12\oint_\gamma(x^\mu\,dx^\nu-x^\nu\,dx^\mu)
      =\oint_\gamma x^\mu\,dx^\nu .
  \]
  Consequently
  \[
    \delta_\gamma\psi
      =i\,t_\alpha F_{\alpha\mu\nu}(X)\psi(X)
        \oint_\gamma x^\mu\,dx^\nu+O(\hbox{area}^{3/2}),
  \]
  and the requested proportionality constant is \(i\).

  Now suppose \(F_{\mu\nu}=0\) in a simply connected finite region
  \(\mathcal R\).  Parallel transport from a fixed point \(x_0\) to \(x\)
  is path independent there: two paths differ by a closed loop, and the
  latter can be tiled by small loops with trivial holonomy.  Let \(U(x,x_0)\)
  be this transport matrix.  It obeys
  \[
    (\partial_\mu-iA_\mu)U(x,x_0)=0.
  \]
  The finite gauge transformation by \(U^{-1}\) therefore changes the
  connection to
  \(A'_\mu=0\) throughout \(\mathcal R\).  The qualification about the
  region is essential: a flat connection can have nontrivial holonomy on a
  multiply connected space.

  \WeinbergSolution{6}
  The quadratic gauge-field kernel in momentum space is
  \[
    K_{\mu\nu}
      =p^2\eta_{\mu\nu}-p_\mu p_\nu+\xi^{-1}n_\mu n_\nu .
  \]
  Direct multiplication shows that its inverse is
  \[
    \Delta_{\alpha\mu,\beta\nu}(p)
    =\frac{\delta_{\alpha\beta}}{p^2-i\epsilon}
      \left[
        \eta_{\mu\nu}
        -\frac{p_\mu n_\nu+n_\mu p_\nu}{p\cdot n}
        +\frac{n^2+\xi p^2}{(p\cdot n)^2}p_\mu p_\nu
      \right].
  \]
  As usual in an axial gauge, the factors \(1/(p\cdot n)\) require a
  specified pole prescription.  In the sharp-gauge limit \(\xi\to0\), the
  propagator is transverse to \(n^\mu\).

  The gauge variation of \(f_\alpha=n^\mu A_{\alpha\mu}\) is
  \[
    \delta f_\alpha=n^\mu(D_\mu\lambda)_\alpha,
  \]
  and hence
  \[
    \mathcal L_{\rm GH}
      =\omega_\alpha^*n^\mu(D_\mu\omega)_\alpha
      =\omega_\alpha^*n^\mu\partial_\mu\omega_\alpha
       +C_{\alpha\gamma\beta}
        \omega_\alpha^*n^\mu A_{\gamma\mu}\omega_\beta .
  \]
  The free ghost propagator is the inverse of \(n\cdot\partial\):
  \[
    \Delta_{\alpha\beta}(p)
      =\frac{\delta_{\alpha\beta}}{i\,n\cdot p}
  \]
  with a compatible prescription.  With all momenta incoming, the
  ghost--antighost--gauge vertex is proportional to
  \(C_{\alpha\gamma\beta}n^\mu\) (and acquires the conventional factor \(i\)
  when the usual \(i\mathcal L_{\rm int}\) Feynman rule is used).  In the
  strict axial gauge \(n\cdot A=0\), this interaction vanishes on the gauge
  slice; the determinant reduces to \(\det(n\cdot\partial)\), so the ghosts
  decouple.

  \WeinbergSolution{7}
  Assign dimensions
  \[
    [A_\mu]=[\omega]=[\omega^*]=1,\qquad [h]=2,
  \]
  and ghost numbers \(0,1,-1,0\), respectively.  The cohomology argument of
  Section~15.8 says that every local BRST-invariant action of ghost number
  zero has the form
  \[
    I=I_0[A]+s\Psi ,
  \]
  where \(I_0\) is gauge invariant and \(\Psi\) is a Lorentz- and
  global-gauge-invariant fermion of ghost number \(-1\).

  For a semisimple gauge algebra, the gauge-field part of the most general
  density of dimension at most four is, up to a constant, a surface term,
  and independent invariant quadratic forms on simple factors,
  \[
    \mathcal L_0
      =-\frac14 Z_{\alpha\beta}
        F_\alpha^{\mu\nu}F_{\beta\mu\nu}
       +\frac14\vartheta_{\alpha\beta}
        F_\alpha^{\mu\nu}\widetilde F_{\beta\mu\nu}.
  \]
  (For Abelian factors one may add the familiar gauge-invariant terms
  permitted by their invariant tensors.)  Modulo integrations by parts, the
  most general renormalizable gauge-fixing fermion density is
  \[
    \begin{split}
    \psi={}&
      z_{\alpha\beta}\omega_\alpha^*\partial_\mu A_\beta^\mu
      +\frac12\xi_{\alpha\beta}\omega_\alpha^*h_\beta
      +a_{\alpha\beta\gamma}
       \omega_\alpha^*A_\beta^\mu A_{\gamma\mu}\\
      &+\frac14 b_{\alpha\beta\gamma}
       \omega_\alpha^*\omega_\beta^*\omega_\gamma ,
    \end{split}
  \]
  where all coefficient tensors are invariant under the global gauge group,
  \(a_{\alpha\beta\gamma}=a_{\alpha\gamma\beta}\), and
  \(b_{\alpha\beta\gamma}=-b_{\beta\alpha\gamma}\).  Thus the answer is
  \[
    \boxed{\mathcal L=\mathcal L_0+s\psi}
  \]
  up to total derivatives and field-independent constants.  The usual
  linear covariant gauge is obtained with \(z_{\alpha\beta}\propto
  \delta_{\alpha\beta}\), \(\xi_{\alpha\beta}\propto\delta_{\alpha\beta}\),
  and \(a=b=0\).  The last two terms display interactions allowed by the
  stated principles but absent in that special Faddeev--Popov choice,
  including, after applying \(s\), nonlinear gauge-fixing and four-ghost
  interactions.

  \WeinbergSolution{8}
  Let \(\epsilon_F\) denote the Grassmann parity of \(F\).  Moving a
  functional derivative through a homogeneous functional and exchanging
  left and right derivatives in Eq.~\eqref{eq:15.9.13} gives
  \[
    (F,G)
      =-(-1)^{(\epsilon_F+1)(\epsilon_G+1)}(G,F).
  \]
  If \(F\) and \(G\) are both bosonic the sign is \(+\); in every other
  case it is \(-\), precisely Eq.~\eqref{eq:15.9.15}.

  For the Jacobi identity, collect fields and antifields into coordinates
  \(z^A=(\chi^n,\chi_n^\ddagger)\).  The bracket can be written
  \[
    (F,G)=F\frac{\overleftarrow\delta}{\delta z^A}
      E^{AB}\frac{\overrightarrow\delta G}{\delta z^B},
  \]
  where \(E^{AB}\) is the constant odd symplectic matrix with only
  field--antifield blocks nonzero.  Expanding
  \((F,(G,H))\), the terms containing one second derivative cancel in the
  graded cyclic sum because \(E^{AB}\) has graded antisymmetry; the terms
  containing a second derivative of \(G\) or \(H\) cancel in the same way.
  Hence
  \[
    (-1)^{(\epsilon_F+1)(\epsilon_H+1)}(F,(G,H))
      +\hbox{cyclic permutations}=0.
  \]
  Its first coefficient is \(-1\) only when \(F\) and \(H\) are both
  bosonic, and is \(+1\) otherwise, which is the sign convention of
  Eq.~\eqref{eq:15.9.21}.

  \WeinbergSolution{9}
  There is a necessary editorial correction in the stated condition.  The
  quantum-observable condition is
  \[
    (O,S)=i\Delta O,
    \qquad\text{or}\qquad
    \sigma O\equiv(O,S)-i\Delta O=0,
  \]
  as also stated in the footnote following
  Eq.~\eqref{eq:15.9.36}.  The printed right-hand side \(i\Delta S\) does
  not imply gauge-fixing independence for a general \(O\).

  To prove the corrected statement, vary the gauge fermion by
  \(\delta\Psi\), so the Lagrangian submanifold
  \(\chi_n^\ddagger=\delta\Psi/\delta\chi^n\) is shifted
  anticanonically.  Field-space integration by parts, including the
  Jacobian measured by \(\Delta\), gives the normalized variation
  \[
    \delta_\Psi\langle O\rangle
      =i\bigl\langle \sigma(O\,\delta\Psi)\bigr\rangle
       -i\langle O\rangle\bigl\langle\sigma(\delta\Psi)\bigr\rangle .
  \]
  BV integration by parts makes the expectation value of every
  \(\sigma\)-exact insertion vanish when
  \[
    \frac12(S,S)-i\Delta S=0.
  \]
  Moreover, the quantum master equation implies \(\sigma^2=0\), and the
  graded product rule reduces the remaining term to one proportional to
  \(\sigma O\).  It therefore vanishes for a quantum observable:
  \[
    \delta_\Psi\langle O\rangle=0.
  \]
  When \(O\) depends only on the fields, \(\Delta O=0\), so the condition
  reduces to ordinary generalized BRST invariance \((O,S)=0\), in agreement
  with Eq.~\eqref{eq:15.9.36}.
\end{enumerate}
